% 附录 A-D
\appendix
\chapter{常见客户问题 FAQ}

\begin{quote}
以下问答全部基于 EdgeFlow v0.1.0（2026-08-14 发布）已核验实现作答；涉及规划中能力均明确标注。
\end{quote}

\section*{A1. 云边断网期间，采集到的数据会丢吗？}
\addcontentsline{toc}{section}{A1. 云边断网期间，采集到的数据会丢吗？}

断网期间边缘本地采集不受影响：设备 $\rightarrow$ Mapper $\rightarrow$ 设备影子的链路完全在边缘本地完成，影子始终保留最新状态快照。但平台不保留历史遥测样本，断网期间产生的中间样本不会回补，重连后通过周期上报（默认 30s，可配 1s--10min）补报\textbf{最新状态快照}。如业务需要完整历史序列，需在 EventBus 侧或应用侧自行订阅持久化（v0.1.0 未内置遥测历史存储）。

\section*{A2. 弱网下发给设备的指令会丢吗？}

不会静默丢失。下行指令走可靠投递：应用层 Ack 确认、同一命令 ID 自动重试（5s 超时 $\times$ 最多 3 次尝试）、边缘侧幂等去重（缓存 1000 条 FIFO），通道恢复后重发指令仍可正确投递且不重复执行；超时最终返回 504 便于调用方感知。

\section*{A3. 数据是否可追溯？}

可追溯。三类台账覆盖全链路：审计台账（云端 API 操作，含 401 拒绝记录，JSONL）、ops-ledger（keadm 升级/回滚，含备份 id 与失败原因）、op-ledger（设备级采集/下发流水，SQLite，30 天保留）。统一口径：时间戳（毫秒/RFC3339）+ 设备标识（nodeID/deviceName/namespace）+ 同步状态（result/status）。

\section*{A4. 支持哪些设备协议？}

v0.1.0 支持 MQTT（QoS1，遥测/指令主题）与 Modbus TCP（显式设置 \env{EDGEFLOW\_MODBUS\_ADDR} 启用，读写保持寄存器/线圈，写后回读验证）；另提供 mock\_sensor 模拟器用于演示联调。OPC-UA 接入\warn。

\section*{A5. 能管理多少边缘节点？}

支持多节点统一纳管。性能基线上 10 节点并发注册实测 100\% 成功（平均 201.3ms）；更大规模请以实际环境压测为准。

\section*{A6. 平台安全机制有哪些？}

生产加固四项：mTLS 云边通道（启用后自动升级 wss）、接入 Token 鉴权（默认关闭、开启后 401 fail-fast）、审计台账（文件 0600/目录 0700，认证失败也记录）、镜像多架构与 keadm 升级备份校验。完整 RBAC（当前为单 Token 鉴权）与镜像安全扫描\warn。

\section*{A7. 模型如何部署到边缘？}

v0.1.0 以“容器镜像 + 配置”为载体：模型与推理运行时打包为版本化镜像 $\rightarrow$ podsync 下发 Pod 声明（image + replicas）$\rightarrow$ 边缘 Edged 拉取运行并持续自治（多副本/健康自愈/防漂移）；推理参数经 config-sync 下发并落盘。模型仓库/版本管理/灰度发布平台\warn，当前以镜像 Tag + 配置参数化作为过渡路径。

\section*{A8. 升级和回滚有保障吗？}

有。keadm upgrade 自动备份（\seqsplit{\texttt{manifest.json}} + sha256 校验）$\rightarrow$ staging 原子替换 $\rightarrow$ 支持 \seqsplit{\texttt{--simulate-failure}} 演练；失败后执行 \seqsplit{\texttt{keadm rollback --latest}} 可恢复，全程 ops-ledger 留痕（operator/from/to/result/note）。配合镜像 Tag 切换可完成应用级升级与回滚。

\section*{A9. 边缘重启或云端重启会怎样？}

边缘重启：期望状态（节点元数据/Pod/配置）已落盘 MetaManager SQLite（WAL），启动后自动恢复并按声明调谐。云端重启：云端状态为内存态，边缘感知断线（读超时约 90s）后退避重连并自动重新注册（首条消息为 Register），随后周期上报自动重建节点/设备/Pod 状态，无需人工介入（E2E 已实测）。

\section*{A10. 设备离线如何感知？}

v0.1.0 提供\textbf{节点级}离线判定（CloudHub 90s 无消息断开 + NodeController 心跳停滞 180s 双路径），状态呈现 Ready/Unknown/Offline；\textbf{设备级}离线检测\warn，当前 Mapper 采集失败仅记 Warn 并落 op-ledger（result=error），可通过台账查询错误历史。

\chapter{术语表}

{\setlength{\tabcolsep}{3pt}
\begin{longtable}{p{3.2cm}p{11.8cm}}
\toprule
术语 & 定义\\
\midrule
\endhead
\bottomrule
\endfoot
CloudCore & EdgeFlow 云端组件，承担节点/设备管控、REST API、审计台账与运行指标\\
EdgeCore & EdgeFlow 边缘组件，由 EdgeHub/Edged/DeviceTwin/EventBus/Mapper/MetaManager 组成\\
EdgeHub & 边缘云边通道客户端，负责注册、心跳、消息收发与幂等去重\\
Edged & 边缘容器自治引擎，按声明式调谐（默认 5s）维护容器期望状态\\
DeviceTwin & 设备影子，维护设备属性最新快照（内存态），支撑上报与指令闭环\\
Mapper & 设备接入框架，提供 DeviceMapper 接口与注册路由，实现协议适配（MQTT/Modbus）\\
EventBus & 边缘 MQTT 数据面，承载遥测/指令主题（QoS1），broker 缺失时降级本地模式\\
MetaManager & 边缘元数据管理，SQLite（WAL）落盘节点/Pod/配置，提供本地持久化\\
设备影子 & 设备属性在平台侧的期望/实际状态视图（desired/reported）\\
声明式调谐 & 以期望状态为基准周期对账（reconcile），持续收敛实际状态\\
调谐周期 & 默认 5s，Edged 每次对账的时间间隔\\
镜像漂移 & 容器实际镜像与期望镜像不一致，触发自动重建还原\\
CrashLoopBackOff & 连续重启 3 次进入 30s 退避窗口，稳定运行 60s 后重置计数\\
指数退避重连 & 断线后 1s 起步、逐次翻倍、封顶 60s 的自动重连策略\\
可靠投递 & 应用层 pending+Ack 确认、同 ID 重发、最多 3 次尝试、幂等去重（1000 条）\\
幂等 & 同一命令/消息重复投递不产生重复副作用\\
DeviceReport & 设备状态周期上报消息（30s，可配），无 Ack 单向流式\\
podsync & 云端向边缘下发 Pod 声明（image/replicas）的 API\\
config-sync & 云端向边缘下发配置并持久化的 API\\
节点离线 & 云端判定：CloudHub 90s 断开或 NodeController 180s 心跳停滞\\
设备离线 & 设备级离线检测（\warn）；当前仅 op-ledger 记录采集错误\\
台账 & 追加式审计/操作记录（审计台账/ops-ledger/op-ledger 三类）\\
JSONL & JSON Lines，每行一条 JSON 记录的追加式日志格式\\
mTLS & 双向 TLS，云边通道启用后自动升级 wss\\
QoS1 & MQTT 服务质量级别“至少一次”，不保证去重，消费方需幂等\\
keadm & EdgeFlow 部署工具（init/join/upgrade/rollback/ops-ledger 等）\\
ops-ledger & keadm 升级/回滚操作台账（JSONL，含备份 id 与失败原因）\\
op-ledger & 设备操作台账（SQLite 追加流水，30 天保留；代码/SQL 表名为 op\_ledger）\\
审计台账 & 云端 API 操作台账（audit-ledger.jsonl，含 401 记录）\\
状态快照周期补报 & 弱网恢复后按周期补报最新状态快照（非历史回补）\\
\end{longtable}}

\chapter{产品特性-场景映射表}

\textbf{表 C-1：F01--F40 已实现特性 × 四场景归属（双向追溯）}

{\setlength{\tabcolsep}{3pt}
\begin{longtable}{p{1.0cm}p{4.6cm}ccccp{3.2cm}}
\toprule
F 编号 & 特性名称 & \makecell{数据\\采集} & \makecell{弱网\\自治} & \makecell{模型\\管理} & \makecell{模型\\应用} & 手册引用章节\\
\midrule
\endhead
\bottomrule
\endfoot
F01 & WebSocket 云边管理面通道 & $\bullet$ & $\bullet$ & $\circ$ & $\circ$ & 1.2 / 3.3\\
F02 & 节点注册（Register/RegisterAck，10s 超时） & $\circ$ & $\bullet$ & $\circ$ & $\circ$ & 3.2 / 3.3\\
F03 & 心跳保活（30s） & $\circ$ & $\bullet$ & $\circ$ & $\circ$ & 3.3 / 3.4\\
F04 & 指数退避重连（1s$\rightarrow$60s 封顶） & $\circ$ & $\bullet$ & $\circ$ & $\circ$ & 3.2 / 3.3\\
F05 & 消息可靠投递（5s×3、幂等去重 1000 条、504） & $\bullet$ & $\bullet$ & $\bullet$ & $\circ$ & 2.3 / 3.2 / 4.3 / 6.1\\
F06 & 多节点并发管理（10 节点实测 100\%） & $\bullet$ & $\circ$ & $\circ$ & $\circ$ & 附录 A5\\
F07 & 在线/离线双路径判定（90s/180s） & $\circ$ & $\bullet$ & $\circ$ & $\circ$ & 3.3 / 3.6\\
F08 & 节点状态语义（Ready/Unknown/Offline） & $\circ$ & $\bullet$ & $\circ$ & $\circ$ & 3.3\\
F09 & 云端状态内存态+边缘持久（自动恢复） & $\circ$ & $\bullet$ & $\circ$ & $\circ$ & 3.3 / 3.6\\
F10 & 消息路由引擎（SendToNode/Broadcast/Deliver） & $\circ$ & $\circ$ & $\circ$ & $\circ$ & 附录 C（仅登记）\\
F11 & Edged 声明式调谐（5s） & $\circ$ & $\bullet$ & $\bullet$ & $\bullet$ & 3.2 / 4.3 / 5.3\\
F12 & 多副本 Pod 管理（补齐/收缩） & $\circ$ & $\bullet$ & $\bullet$ & $\bullet$ & 3.3 / 4.2 / 5.3\\
F13 & 健康自愈自动重启 & $\circ$ & $\bullet$ & $\bullet$ & $\bullet$ & 3.3 / 4.3 / 5.3\\
F14 & CrashLoopBackOff（3 次/30s/60s 重置） & $\circ$ & $\bullet$ & $\bullet$ & $\bullet$ & 3.3 / 4.3 / 5.3\\
F15 & 镜像漂移检测与重建 & $\circ$ & $\bullet$ & $\bullet$ & $\bullet$ & 3.3 / 4.3 / 5.3\\
F16 & 断网自治（E2E 60s 实测） & $\circ$ & $\bullet$ & $\circ$ & $\bullet$ & 3.3 / 3.6 / 5.3\\
F17 & MetaManager SQLite 持久化 & $\circ$ & $\bullet$ & $\bullet$ & $\bullet$ & 3.3 / 4.3 / 5.3 / 6.1\\
F18 & PodStatus 周期上报（30s 可配，Absent 90s） & $\circ$ & $\bullet$ & $\circ$ & $\circ$ & 3.3 / 3.4\\
F19 & 删除收敛闭环（delete$\rightarrow$Absent$\rightarrow$云端收敛） & $\circ$ & $\bullet$ & $\circ$ & $\circ$ & 附录 C（仅登记）\\
F20 & MetaManager 增量订阅通知 & $\circ$ & $\bullet$ & $\circ$ & $\circ$ & 3.3\\
F21 & DeviceTwin 影子（desired/reported 合并） & $\bullet$ & $\circ$ & $\circ$ & $\bullet$ & 2.3 / 5.3\\
F22 & 设备属性周期上报（30s 可配，启动即报） & $\bullet$ & $\circ$ & $\circ$ & $\bullet$ & 2.3 / 5.3 / 6.1\\
F23 & device-command 指令闭环 & $\bullet$ & $\circ$ & $\circ$ & $\bullet$ & 2.3 / 5.3 / 6.1\\
F24 & Mapper 框架（接口+注册路由） & $\bullet$ & $\circ$ & $\circ$ & $\circ$ & 2.2 / 2.3\\
F25 & mock\_sensor 模拟器（2s，targetTemp/reset） & $\bullet$ & $\circ$ & $\circ$ & $\circ$ & 2.3 / 2.4\\
F26 & Modbus TCP Mapper（显式启用，写后回读） & $\bullet$ & $\circ$ & $\circ$ & $\circ$ & 2.3 / 2.4\\
F27 & op-ledger 操作台账（SQLite 表名 op\_ledger，30 天保留） & $\bullet$ & $\circ$ & $\circ$ & $\circ$ & 2.3 / 2.5 / 6.3\\
F28 & EventBus MQTT 数据面（QoS1，自动重连） & $\bullet$ & $\circ$ & $\circ$ & $\bullet$ & 2.3 / 5.3\\
F29 & EventBus 降级本地模式 & $\bullet$ & $\bullet$ & $\circ$ & $\circ$ & 2.3 / 2.4\\
F30 & 设备 API 云端视图 & $\bullet$ & $\circ$ & $\circ$ & $\bullet$ & 2.3 / 5.3 / 6.1\\
F31 & REST API 13 端点（podsync/config-sync 等） & $\circ$ & $\circ$ & $\bullet$ & $\bullet$ & 4.3 / 5.4\\
F32 & 运行指标 5 项（edgeflow\_cloudcore\_*） & $\circ$ & $\circ$ & $\circ$ & $\circ$ & 1.2 / 1.4\\
F33 & Token 认证（默认 off，401 fail-fast） & $\circ$ & $\circ$ & $\circ$ & $\circ$ & 附录 A6\\
F34 & mTLS 云边通道（wss） & $\circ$ & $\circ$ & $\circ$ & $\circ$ & 附录 A6 / 3.3\\
F35 & 审计台账（JSONL，0600/0700） & $\circ$ & $\circ$ & $\circ$ & $\circ$ & 6.3\\
F36 & keadm 部署（init/join/ops-ledger） & $\circ$ & $\circ$ & $\bullet$ & $\circ$ & 4.3\\
F37 & Helm Chart 部署 & $\circ$ & $\circ$ & $\bullet$ & $\circ$ & 4.3\\
F38 & 多架构镜像（amd64+arm64） & $\circ$ & $\circ$ & $\bullet$ & $\circ$ & 4.3\\
F39 & 升级回滚机制（备份/演练/审计/留痕） & $\circ$ & $\circ$ & $\bullet$ & $\circ$ & 4.3 / 4.4\\
F40 & 性能基线（10 节点 201.3ms） & $\circ$ & $\circ$ & $\circ$ & $\circ$ & 附录 A5\\
\end{longtable}}

\begin{quote}
图例：$\bullet$ 主要支撑场景；$\circ$ 次要/间接支撑场景。手册引用章节为 v1.0.0 定稿节号。
\end{quote}

\textbf{表 C-2：规划中/范围外特性清单（F41--F48）}

{\small\setlength{\tabcolsep}{3pt}
\begin{tabular}{@{}p{1.0cm}p{3.4cm}p{2.6cm}p{7.6cm}@{}}
\toprule
编号 & 规划项 & 状态 & 说明\\
\midrule
F41 & 模型仓库与版本管理平台 & \warn & 当前以镜像 Tag + 配置参数化过渡（第 4 章）\\
F42 & 灰度发布 & \warn & ROADMAP 已规划；过渡期人工分批（4.2）\\
F43 & 完整 RBAC（多角色权限） & 规划中 & 当前为单 Token 鉴权\\
F44 & 设备认证（设备级身份） & 规划中 & 设备级离线检测同属规划范围\\
F45 & 调度/超卖、Flannel 网络 & 规划中 & 集群化能力增强\\
F46 & OPC-UA 协议接入 & 规划中 & 当前支持 MQTT/Modbus TCP\\
F47 & 通道压缩/Protobuf 编码 & 规划中 & 当前 JSON 信封\\
F48 & 镜像安全扫描 & 规划中 & 当前提供 sha256 完整性校验（升级备份）\\
\bottomrule
\end{tabular}}

\chapter{版本管理与更新说明（回滚口径）}

\begin{enumerate}
  \item \textbf{本手册版本管理}：以 Git 提交为版本锚点，主源文件为 Markdown（本文件），HTML/LaTeX/PDF 为同版本衍生制品；修订须更新“修订记录”表并重新生成全部衍生制品。
  \item \textbf{更新说明}：产品新版本发布后，本手册须对照特性清单（附录 C）逐条核验“已实现/规划中”状态后再升版。
  \item \textbf{回滚口径}：若某版本手册被发现事实性错误，回退到上一 Git 提交并发布勘误版本（vX.Y.Z+1），不直接修改已发布历史版本文件；勘误记录写入修订记录表。
  \item \textbf{评审留痕}：v1.0.0 由方案团队（主笔）+ 产品团队（事实基线提供）+ 独立技术复核（特性$\leftrightarrow$手册双向追溯、结构完整性、异常路径、术语一致性）完成，复核记录见《评审记录与交接说明.md》。
\end{enumerate}
